% =============================================================================
%  扉页、译序、读法与体例
% =============================================================================
\begin{titlepage}
\vspace*{1.2cm}
\begin{center}
{\bfseries\LARGE 通过膜探测量子化}\\[0.35cm]
{\large\itshape Probing Quantization Via Branes}\\[1.2cm]
{\large Davide Gaiotto}\\[0.1cm]
{\small\itshape Perimeter Institute, Waterloo, Ontario, Canada N2L 2Y5}\\[0.45cm]
{\large Edward Witten}\\[0.1cm]
{\small\itshape Institute for Advanced Study, Einstein Drive, Princeton, NJ 08540 USA}\\[1.6cm]
{\large 中文译注本}\\[0.5cm]
{\small 底本：arXiv:2107.12251v2（2021 年 8 月 11 日），共 101 页}\\[0.15cm]
{\small 原文标题：Probing Quantization Via Branes}\\[0.15cm]
{\small 原文日期：2020 年 5 月}\\[1.2cm]
{\large 数译（Shuyi）\ 译注}\\[0.4cm]
{\small 2026 年 11 月}
\end{center}
\vfill
\begin{center}
\begin{minipage}{0.86\textwidth}
\small
\textbf{权利说明。}原文版权归原作者所有。本译本为非营利性的学术翻译，供学习与研究参考；译文、译注与排版由数译项目完成。引用本文内容时请以原文为准，并注明原文出处。若权利人认为本译本有不妥之处，请与译者联系。
\end{minipage}
\end{center}
\end{titlepage}

\tableofcontents
\clearpage

\section*{译序}
\addcontentsline{toc}{section}{译序}

Gaiotto 与 Witten 的《通过膜探测量子化》写于 2020 年 5 月，2021 年 8 月修订后挂出。它要做的事情可以用一句话概括：把“量子化”这件事放进二维拓扑场论的语境里重讲一遍，然后回答一个很自然的问题——用膜做出来的量子化，和用极化做出来的几何量子化，什么时候是同一件事？

这个问题并不小。几何量子化自 Kostant、Souriau 以来已相当成熟，但它有一个众所周知的软肋：结果一般依赖极化的选取；而不同极化给出等价理论的情形，往往正是有意思的物理。膜量子化（Gukov–Witten, 2008）换了一个起点：不选极化，而选辛流形 $M$ 的一个复化 $Y$，把 $M$ 的量子 Hilbert 空间定义为 $A$-模型中 $(\B,\B_\cc)$ 弦的物理态空间。两种方法都依赖“额外结构”的选取，于是“它们何时一致”就变成了“极化和复化何时相容”的问题。

本文的答案是：当极化 $\P$ 可以解析延拓为 $Y$ 的全纯极化时，两者一致（第 3 节）。围绕这条主线，文章还讨论了对称性（第 4 节）与复辛流形作为实辛流形的量子化（第 5 节），并在附录里补齐了反常、$\Hom$ 的微扰计算、相反代数等细节。

这篇文字不好读，原因有三。其一，它横跨几何量子化、形变量子化、$A$-模型与 $A$/B-膜、超对称 $\sigma$-模型等多个领域，默认读者对每一块都有相当的准备。其二，作者大量使用“物理上显然、数学上待补”的论证，很多关键步骤只给图景不给证明——第 3、4 节的若干结论至今仍应视为纲领性的。其三，原文自身的记号相当紧缩（单字母宏、复化与实化混用、$\B$ 与 $\B_\cc$ 的角色随语境切换），稍不留神就会读错方向。

因此本译本在忠实直译之外，做了三件事：术语首现处附英文原词；对理解主线必需的背景概念补“概念卡片”；对原文的笔误、记号冲突与论证跳步，用“译注”与“校勘”标出，并尽量给出可核查的依据。凡属译者补出的解释，一律标注“译者所加”或“无外部依据”，不与原文内容混同。

翻译所据底本为 arXiv:2107.12251v2 的 PDF，同时使用了作者提供的 LaTeX 源码（Draft\_D.tex）。公式、编号、参考文献均沿用原文，以便与原文逐页对照。

\section*{读法}
\addcontentsline{toc}{section}{读法}

\begin{itemize}[leftmargin=1.6em,itemsep=0.25em]
\item \textbf{只想抓主线}：读第 1 节、第 3.1 节、第 3.7--3.8 节、第 4.1 节与第 4.5--4.6 节，再回看第 2.5--2.7 节的定义。
\item \textbf{关心几何量子化}：第 2.2 节（极化）、第 3.7 节（实极化）、第 3.8 节（复极化）是核心；第 3.10 节给出平坦 $G$-丛模空间这个最重要的例子。
\item \textbf{关心对称性与几何 Langlands}：第 4 节整节，特别是 4.4--4.6；第 5 节是其在复辛流形上的应用。
\item \textbf{关心技术细节}：附录 A（反常的显式计算）、附录 B（$\Hom$ 的最低阶与修正）、附录 C（相反代数）、附录 D（$\B_\cc$ 的反常）。
\end{itemize}

原文页码与 PDF 页码相差 2（封面、目录各一页）。正文中“原文第 $n$ 页”指印刷页码。

\section*{体例}
\addcontentsline{toc}{section}{体例}

\begin{itemize}[leftmargin=1.6em,itemsep=0.2em]
\item 术语首次出现时写作“中文（English）”；此后一般只用中文。关键术语的对照见文末《术语对照表》。
\item \textbf{译注}：解释背景、补足跳步、说明译法取舍。凡引用外部材料，标注来源等级：A 级为原始论文或权威专著，B 级为 nLab、Stacks Project、ProofWiki 等经交叉核对的条目，I 级为项目内部术语库；无外部依据者标“译者按（无外部依据）”。
\item \textbf{校勘}：标出原文的笔误、记号冲突或疑似错误，给出原文写法、建议改法与依据。
\item \textbf{概念卡片}与\textbf{本节重点}为译者所加，便于初读，不属原文。
\item 公式编号、定理编号、参考文献编号沿用原文；图号按本译本顺序重排，图题译出。
\item 原文斜体强调在中译本中改为粗体（原稿把 \verb|\emph| 重定义为无衬线粗体）。
\end{itemize}
